\chapter{El código que el municipio se dio a sí mismo}

\label{cap:cot}
\section{El texto, leído}

El capítulo \ref{cap:metodo} consigna que el texto operativo de este Código --- ciento cincuenta y dos artículos, ocho anexos y un atlas de ocho planos --- fue obtenido y leído. Un barrido posterior, hecho sobre el ejemplar publicado en el sitio del municipio con el criterio de leer por nombre cada mención territorial, precisó dos cosas que el relevamiento inicial no había extraído.

\subsection*{Las dos márgenes del río son no urbanizables}

\begin{ficha}{Código de Ordenamiento Territorial, Título 3, artículo 43}
«VI. \textbf{Zonas PU}: Las zonas establecidas como Parque Urbano, corresponden a sectores
\textbf{no urbanizables} destinados para uso recreativo, la protección del medio ambiente y el
paisaje. Representan los \textbf{ejes verdes urbanos de ribera, hacia ambas márgenes del Río
La Caldera} y componen el sistema de infraestructura verde del municipio generando corredores
ecológicos, longitudinales y transversales según las escorrentías naturales, de conexión hacia
el área de amortiguación.»

El Anexo VI reitera, en la planilla de esa zona: «Costanera / ribera del Rio La Caldera», con
uso dominante «Parque urbano, zona no urbanizable».
\end{ficha}

\textbf{Ambas márgenes del río están declaradas no urbanizables por zonificación municipal.}
Es una protección más amplia que la línea de ribera provincial: no depende de una
determinación caso por caso, rige por zona y es anterior a cualquier expediente.

\textbf{Pero el código no dice cuántos metros.} En ninguna parte del texto hay un ancho: no hay
metros, no hay cota, no hay línea de ribera determinada. \textbf{El ancho de esa franja sólo existe gráficamente}, en el plano del Anexo II que integra el atlas del Código. Este libro no pudo leerlo: en el ejemplar disponible el trazado no resulta mensurable, y el Código no lo traduce a metros en ninguna parte de su articulado.

\begin{cautela}
Eso es, exactamente, lo que el capítulo \ref{cap:opacidad} describe. La norma enuncia una protección fuerte ---
«no urbanizable», «ambas márgenes» --- y \textbf{remite a un plano la única magnitud que
permitiría aplicarla}. Sin esa magnitud no puede saberse si una parcela está dentro o fuera, ni
verificarse ningún loteo contra la zonificación de su propio municipio.

\textbf{Obtener ese plano en escala legible --- o que el municipio informe el ancho en metros --- era, hasta hace poco, la pieza faltante más importante del capítulo \ref{cap:loteo}.} El apartado siguiente cuenta cómo se resolvió, y la respuesta no es la que este libro esperaba.
\end{cautela}

\subsection*{Los dos planos aparecen, y lo que aparece no alcanza}

Ese reclamo queda parcialmente satisfecho, y por una vía que este relevamiento no había
previsto. Los planos del atlas se publican, dentro del PDF del Código, sólo como imágenes a escala municipal y sin las capas con que se confeccionaron; y el propio municipio los reprodujo además dentro del informe que las demandadas elevaron al juzgado en la ejecución de sentencia del
amparo que analiza el capítulo \ref{cap:amparo}. Allí están el \textbf{Anexo I, «Clasificación del
Territorio y Límites»} (lámina \ref{lam:anexo1}) y el \textbf{Anexo II, «Plano de Zonificación
de Usos del Suelo»}, que la lámina \ref{lam:anexo2} reproduce en la versión en color del PDF municipal.

\textbf{El Anexo I clasifica el territorio en tres categorías y traza cuatro límites.} La
leyenda distingue \textbf{A.N.} (área natural), \textbf{A.R.} (área rural) y \textbf{A.U.}
(área urbana), y agrega ejido urbano (E.U.), límite jurisdiccional, línea de servicios (L.S.)
y núcleos urbanos (N.U.). El área urbana es una franja angosta sobre el valle del río; el
resto del municipio se reparte entre rural y natural. Es la primera vez que este relevamiento
ve dibujada la división que el Código enuncia en su articulado.

\textbf{El Anexo II zonifica los usos, y llega acompañado de algo que el plano solo no da:} la
lista de las categorías que habilitan urbanizar. El informe la transcribe ---
\textbf{PA} (protección activa), \textbf{NR} (núcleo rural), \textbf{M1} (mixta residencial y
productiva), \textbf{UD} (urbanización diferida), \textbf{R2} (residencial densidad muy baja) y
\textbf{R1} (predominante residencial densidad baja) --- y explica el circuito: quien quiere
lotear pide al municipio la factibilidad de localización, y allí entran en juego esas
categorías.

\begin{cautela}
Cuatro reservas.

\textbf{Primera.} Las grandes manchas rotuladas en el plano del Anexo II --- \textsc{rn},
\textsc{rm}, \textsc{rm 1} --- no coinciden con los códigos de la lista que lo acompaña.
Pueden ser la misma nomenclatura escrita de otro modo o pueden ser dos capas distintas; a esta
resolución no puede establecerse, y el cotejo debe hacerse contra el atlas original.

\textbf{Segunda.} El párrafo que introduce la lista remite a «(Ver Mapa XX)». La referencia
cruzada quedó sin completar en el documento presentado al juzgado.

\textbf{Tercera. El Anexo I aparece dos veces, en dos versiones distintas.} Una hoja lo reproduce con
barra gráfica de diez kilómetros y otra con barra gráfica de veinte, con encuadre diferente.
Este libro no puede establecer cuál corresponde al atlas del Código y cuál fue rearmada para el
informe; reproduce la de veinte, que es la que trae el título y la leyenda completos.

\textbf{Cuarta, y es el punto:} \textbf{ninguno de los dos planos permite medir}. Son capturas
sobre base de Google Earth ---el Anexo I con barra gráfica de diez o de veinte kilómetros según la versión, el Anexo II con escala gráfica de cuatro---, sin cuadrícula de
coordenadas y con la leyenda ilegible en el ejemplar incorporado al expediente. \textbf{El
ancho de la franja no urbanizable de ambas márgenes sigue sin conocerse en metros.} Lo que
cambió no es que el dato aparezca: es que ahora se sabe dónde está dibujado, en qué soporte y
con qué precisión.
\end{cautela}

\textbf{Y hay una asimetría de soporte que conviene retener}, porque vuelve en el capítulo \ref{cap:opacidad}.
La cartografía que la Provincia produjo sobre esta misma cuenca --- los tres mapas temáticos
de mayo de 2022, láminas \ref{lam:mapa1}, \ref{lam:mapa2} y \ref{lam:mapa3} --- está en
\textbf{POSGAR 2007, Faja 3}, a escalas 1:125.000 y 1:40.000, y es superponible con cualquier
plano catastral. La del municipio son capturas de pantalla. Los dos instrumentos ordenan el
mismo territorio, y \textbf{solo uno de los dos puede verificarse contra el catastro}.

\lamina{lam-anexo1-clasificacion.jpg}{lam:anexo1}%
{Anexo I del Código de Ordenamiento Territorial: clasificación del territorio y límites}%
{\textbf{Anexo I. Clasificación del Territorio y Límites}, Municipio de La Caldera. Distingue
área natural (A.N.), área rural (A.R.) y área urbana (A.U.); la leyenda agrega ejido urbano,
límite jurisdiccional, línea de servicios y núcleos urbanos. Sobre base de Google Earth y con
barra gráfica de veinte kilómetros. Encabeza la hoja el título «Zona urbana y no urbana
comprendidas en la micro cuenca», que abrevia la manda de la sentencia del Juzgado de Minas:
determinar «las áreas de zona urbana y no urbana comprendidas en la microcuenca»
(capítulo \ref{cap:amparo}).}%
{Informe de las demandadas en la ejecución de sentencia del amparo colectivo, expediente
INC-23860/1, hoja de zonificación. Reproducción facsimilar; se recortó el margen en blanco del
escaneo y no se alteraron los tonos.}

\lamina{lam-anexo2-zonificacion.jpg}{lam:anexo2}%
{Anexo II del Código de Ordenamiento Territorial: zonificación de usos del suelo}%
{\textbf{Anexo II. Plano de Zonificación de Usos del Suelo}, Municipio de La Caldera. Las
manchas \textsc{pu} ---parque urbano de ribera--- son la franja que el artículo 43 declara no
urbanizable sobre ambas márgenes del río. La cartela declara una escala gráfica de cuatro
kilómetros: \textbf{a esa escala la franja no es mensurable}, y el Código no traduce su ancho
a metros en ninguna otra parte.}%
{Código de Ordenamiento Territorial del Municipio de La Caldera, página 101, en el ejemplar
publicado por el municipio. Reproducción del plano en color; se recortó el encabezado y el pie
de página y no se alteraron los tonos.}

\subsection*{El ancho no está en ninguna parte, y eso es el hallazgo}

El Código completo se obtuvo del sitio del municipio: nueve títulos, ciento cincuenta y dos
artículos, ocho anexos y un atlas de ocho planos. Con el texto íntegro a la vista, la pregunta
que este capítulo venía arrastrando deja de estar abierta, pero no del modo previsto.

\textbf{Se revisaron los tres lugares donde el ancho podría estar. No está en ninguno.}

\textbf{Primero, el articulado.} El artículo 43 define las zonas PU y no trae cifra alguna.
Fuera de él, la palabra «ribera» aparece en el artículo 21 --- que define \textsc{pu} como
«parque urbano de ribera» ---, en el artículo 85 --- accesibilidad, reforestación y destino a
uso público de los predios fiscales ribereños --- y en dos definiciones del glosario del
Anexo VIII. \textbf{En ninguna hay una medida.}

\textbf{Segundo, la planilla de indicadores.} El Anexo VI trae una planilla por zona, y la de
la zona PU es ésta.

\begin{ficha}{Anexo VI, planilla de indicadores --- Zona PU, página 122 del Código}
\textbf{Característica espacial:} costanera / ribera del río La Caldera, zona del A.U.
\textbf{Uso dominante:} parque urbano, zona no urbanizable.

\textbf{Subdivisión mínima --- frente: (en blanco). Fondo: (en blanco).} Superficie:
2.500 m\textsuperscript{2}.

\textbf{Altura máxima: 7 m.} Edificaciones: sólo para habitación o vigilancia.
\textbf{Veredas: 1,20 m de ancho mínimo.}

Condiciones ambientales sobre nivel medio: polvo 4 g/m\textsuperscript{2} cada diez días;
ruido 30 decibeles; humo 0--0 en escala Ringelmann; albedo 20\,\%.
\end{ficha}

\textbf{La planilla fija en metros la altura de lo que se puede construir y el ancho de una
vereda, y deja en blanco el frente y el fondo mínimos.} Del ancho de la franja, nada.

\textbf{Y tercero, el plano.} La cartela del Anexo II declara «esc.\ 4 km»: es una escala
gráfica con la que el municipio entero entra en una página. \textbf{A esa escala, una franja
de decenas de metros ocupa menos de un milímetro de dibujo.} No es que la reproducción sea
mala --- la lámina \ref{lam:anexo2} es ahora el plano oficial en color, tomado del propio PDF
del municipio ---: es que el dibujo no tiene resolución para contener el dato.

\begin{cautela}
Conviene decir con precisión qué se estableció y qué no.

\textbf{Lo establecido:} el ancho de la franja de Parque Urbano \textbf{no existe como cifra
en el Código}. Ni en el articulado, ni en la planilla de indicadores, ni de manera medible en
el plano. Este libro lo buscó en los tres lugares donde podía estar.

\textbf{Lo que no se sigue de eso:} que la zona no exista. Existe, está dibujada sobre ambas
márgenes y tiene contorno definido. Lo que no tiene es traducción a metros dentro del
instrumento.

\textbf{Y hay una vía que el propio Código abre.} Su introducción declara que «\textbf{todos
los planos se confeccionaron en base a capas georreferenciadas oficiales en formato digital},
con el objetivo de facilitar uso y distribución de la información». \textbf{Si esa capa
existe, el ancho es calculable}: no hay que pedir un número, hay que pedir el archivo.

Eso reformula el pedido de este libro. Durante todo el relevamiento se preguntó cuántos metros
mide la franja. \textbf{La pregunta correcta es otra: dónde está el \textit{shapefile}.} Y
tiene una consecuencia adicional: esa capa y las líneas de ribera de la lámina
\ref{lam:mapa2} están en el mismo territorio y admiten superposición. Con las dos se puede
establecer, parcela por parcela, qué loteos caen dentro de la zona no urbanizable de su propio
municipio. \textbf{Es la comprobación más importante que le queda pendiente a este libro, y ya
no depende de conseguir una cifra que no existe.}
\end{cautela}

\subsection*{Y a pocos kilómetros de ahí, el mismo Estado fija un retiro en metros}

Todo lo anterior describe un vacío: una franja declarada no urbanizable cuyo ancho no está
en ninguna parte en metros. Conviene cerrarlo mostrando qué hace el mismo aparato
administrativo cuando sí quiere fijar un retiro.

El capítulo \ref{cap:tierrafiscal} reconstruye una prescripción adquisitiva sobre una parcela
de tres mil doscientos metros cuadrados en el paraje El Gallinato, en este mismo
departamento. Su plano de mensura --- lámina \ref{lam:plano962} --- lleva estampado, abajo a
la izquierda, el recuadro que reproduce la lámina \ref{lam:zonacamino}.

Ahí, para un camino provincial, hay \textbf{dos cifras legibles, una tercera que no lo es y una prohibición}: el ancho de zona
de camino declarado o proyectado, \textbf{14,70 metros}; el ancho que fija la ley, sesenta;
y un sello de la Dirección de Vialidad que pide no ejecutar construcciones permanentes a
menos de cierta distancia del eje de la calzada. Está fechado, firmado y adherido al plano de
la parcela, de modo que cualquiera que consulte ese inmueble lo encuentra.

\begin{cautela}
La comparación es entre dos instrumentos del mismo Estado sobre el mismo territorio, y
conviene formularla sin exceso.

\textbf{Para el camino} hay número, norma citada, fecha, organismo y un soporte ---el plano
de mensura--- que acompaña al inmueble y se consulta al comprarlo.

\textbf{Para ambas márgenes del río}, declaradas no urbanizables por el artículo 43 de este
Código, hay un trazo en un plano a escala de cuatro kilómetros, sin cifra asociada y sin nada que lo
adhiera a las parcelas afectadas.

Este libro no infiere de ahí ninguna intención, y no puede: son organismos distintos, épocas
distintas y regímenes distintos. \textbf{Lo que registra es una capacidad}. El aparato
administrativo provincial sabe fijar retiros en metros, estamparlos sobre un plano y hacerlos
oponibles; lo hace de rutina para una ruta provincial secundaria. Que la franja ribereña del ordenamiento
municipal no tenga equivalente no se explica, entonces, por imposibilidad técnica.

\textbf{Y hay una diferencia de escala que hace la comparación más incómoda, no menos:} el
camino protegido con esas cifras es una ruta provincial; el curso que queda sin cifras es el
de la cuenca cuyos desbordes, y los de sus arroyos, se llevaron casas en 2018 y 2019.
\end{cautela}

\lamina{lam-zona-camino.jpg}{lam:zonacamino}%
{Recuadro de zona de camino y sello de retiro de edificación, Plano 000962}%
{\textbf{Detalle del Plano Nº 000962.} A la izquierda, el recuadro manuscrito de la
Dirección de Vialidad de Salta, del 28 de noviembre de 2014: ruta Provincial Nº 11, tramo 2º
El Desmonte --- Campo Santo, ancho de zona de camino declarado o proyectado 14,70 metros, y
sesenta metros según el artículo 22 inciso a) de la Ley 3.383. A la derecha, el sello que
pide no ejecutar construcciones de carácter permanente a menos de la distancia indicada
--- manuscrita y no legible con certeza en esta reproducción --- medida desde el eje de la
calzada existente.}%
{Anexo del Decreto Nº 1968 del 8 de junio de 2015, Boletín Oficial Nº 19.561. Reproducción
facsimilar ampliada; la hoja se giró noventa grados y se recortó el detalle, sin alterar los
tonos.}

\subsection*{Los límites del municipio, citados de dos maneras}

El artículo 3 fija el ámbito de aplicación: el municipio de La Caldera es el departamento homónimo «exceptuando el territorio correspondiente al municipio de vaqueros», según las leyes provinciales \textbf{2131/47 y 4987/74}.

El artículo 36, diecinueve páginas después, define el mismo límite jurisdiccional invocando tres leyes: \textbf{853/47, 4365/70 y 4987/74}.

\begin{cautela}
\textbf{Corresponde una rectificación.} Una versión anterior de este apartado sostenía que las dos enumeraciones no coincidían y que el Código citaba dos fundamentos distintos para su propio límite. \textbf{No es así.} La Ley 2131 es la 853 renumerada ---el propio Boletín Oficial la publica como «Ley Nº 2131 (Original 853)», la norma que en 1947 fijó la delimitación de los distritos municipales de la provincia; para este departamento dice, en una sola línea, que «el Municipio de la Caldera abarca todo el territorio del Departamento de su mismo nombre», porque entonces no había otro---, y la Ley 4987 de 1974 modificó expresamente la 4365 de 1970. Las dos enumeraciones remiten al mismo régimen: el artículo 3 lo cita por sus leyes vigentes, y el 36 agrega la que la última modificó. Se deja la corrección visible porque el error tiene la forma que la Advertencia preliminar describe: \textbf{una doble numeración leída como dos normas} produjo la apariencia de una contradicción donde no la había.
\end{cautela}

\begin{ficha}{Leyes 4365/70 y 4987/74 --- los límites de Vaqueros}
\textbf{Ley 4365}, sancionada el 5 de agosto de 1970 y publicada en el Boletín Oficial Nº 8.612 del 18 de agosto, bajo el gobierno de facto. Crea «el municipio de \textbf{tercera categoría (2da.\ clase)} de Vaqueros» y le fija estos límites: «al Norte, municipio de La Caldera, \textbf{separado por el río Wierna}; al Sud, departamento de la Capital, \textbf{separado por el río Vaqueros}; al Este, municipio de La Caldera, \textbf{separado por el río La Caldera} y al Oeste finca Lesser y Abra de Lesser, del municipio de La Caldera».

\textbf{Ley 4987}, sancionada el 19 de septiembre de 1974, promulgada el 17 de octubre por el gobernador Ragone y publicada en el Boletín Oficial Nº 9.613 del 30 de octubre. Su artículo 1º establece como límite jurisdiccional entre los dos municipios «\textbf{la línea que forma el río Potrero de Castilla}, desde su nacimiento en el Nevado de Castilla y que continúa en dirección Sudeste para tomar el nombre de \textbf{río Los Yacones} y concluir como \textbf{río Wierna} en su confluencia con el río La Caldera». El artículo 2º modifica en ese punto las leyes 853/47 y 4365/70. El artículo 3º ordena a la Dirección General de Inmuebles \textbf{el deslinde y amojonamiento} del límite en un plazo de noventa días.
\end{ficha}

\textbf{El límite entre La Caldera y Vaqueros corre, entonces, por el agua.} Al norte de Vaqueros es el sistema Potrero de Castilla--Los Yacones--Wierna, hasta la confluencia con el río La Caldera; al este, según la ley de 1970, el propio río La Caldera, un tramo que la ley de 1974 no menciona y que por lo tanto no parece haber modificado. El río Vaqueros es otra cosa: separa a Vaqueros de la capital, y es el límite que el capítulo \ref{cap:vaqueros} documenta como cauce que ninguna de las dos jurisdicciones administra entero. \textbf{Aquí el problema se repite entre los dos municipios del departamento}: cada uno tiene media caja del Wierna, y del río La Caldera aguas abajo de la confluencia.

Tres cosas se siguen de las dos leyes. La primera: el deslinde y amojonamiento de 1974, si se hizo, dejó una mensura en la Dirección General de Inmuebles, y el apéndice \ref{cap:pedidos} la pide. La segunda: la ley que lo ordenó la promulgó el gobernador Ragone, el mismo que un año antes había firmado el Decreto 1.410/73 de loteos (capítulo \ref{cap:loteo}). Y la tercera: la ley de 1970 ubicaba las fincas Lesser y Abra de Lesser «del municipio de La Caldera», mientras el Código de Vaqueros, en su texto de 2026, incluye a Lesser en su ejido (capítulo \ref{cap:vaqueros}); cómo se produjo ese cambio no consta en las dos leyes.

\subsection*{Y hay una versión del límite que sí está dibujada}

Las dos leyes describen el límite con palabras. Existe, además, un trazado dibujado: el que la provincia publica (lámina \ref{lam:idesalim}).

\begin{ficha}{Visor de IDESA, capa «Límites Municipales», 15 de septiembre de 2026}
Bajo la categoría \textbf{Límites y unidades territoriales}, el visor de la Infraestructura de
Datos Espaciales de la provincia publica la capa \textbf{«Límites Municipales»}, junto con
«Departamento», «Gobierno Local» y «Área protegida».

La exportación del visor del 16 de septiembre, que trae la leyenda (lámina \ref{lam:idesamun}), identifica esa capa con \textbf{trazo punteado verde}; el departamento va en punteado oliva. Con esa clave, en la captura ampliada del día anterior el límite municipal \textbf{sigue el río Wierna} desde Los Yacones hasta su confluencia con el río La Caldera, y desde allí continúa hacia el sur. A la altura del pueblo no hay trazo verde: la línea magenta que acompaña allí al río La Caldera pertenece a otra capa, que la captura no identifica. En la vista de conjunto, el límite recorre el oeste del departamento en dirección noroeste--sudeste y alcanza el corredor de la Ruta 9 unos kilómetros al sur del pueblo, que queda entero del lado de su municipio.
\end{ficha}

\begin{cautela}
\textbf{Lo que esto agrega, y es poco pero es concreto.} El visor no declara sobre qué norma se
trazó ese límite, y este relevamiento no lo consultó por vía formal. \textbf{Y el trazado coincide con la Ley 4987.} Leída con la clave de la leyenda, la capa no corre por el río La Caldera aguas arriba de la confluencia: sigue el Wierna, como manda la ley de 1974, y toma el río La Caldera sólo aguas abajo, como la de 1970. Una versión anterior de este apartado leyó como límite municipal la línea magenta que acompaña al río a la altura del pueblo, y dejaba abierta una posible contradicción con la ley; la leyenda la descarta. Por lo mismo, nada en la capa impide que la zona PU del artículo 43 abarque «ambas márgenes» del río en el pueblo. Lo que el visor sigue sin permitir decir es si el municipio reconoce ese trazado, y con qué precisión se dibujó: a la escala de la exportación, un milímetro de pantalla equivale a cientos de metros.

\textbf{Lo que sí se puede decir} es que existe un trazado publicado, que es verificable por
cualquiera y que \textbf{convierte una pregunta abstracta en una comparación}: basta poner ese
límite al lado del que el municipio aplique, y ver si coinciden.

\textbf{Y una observación que conecta con el capítulo \ref{cap:vaqueros}:} el límite entre los dos municipios de esta cuenca \textbf{pasa por el medio del agua}, como el que separa a Vaqueros de la capital. Cada uno tiene media caja del río, y ninguno el cauce entero.
\end{cautela}

\lamina{lam-idesa-limite.jpg}{lam:idesalim}%
{El límite municipal publicado, siguiendo el río Wierna}%
{\textbf{Visor de IDESA, vista ampliada de la confluencia.} En trazo punteado verde, la capa «Límites Municipales» (clave en la lámina \ref{lam:idesamun}): baja desde Los Yacones por el río Wierna, rotulado a la izquierda, y continúa hacia el sur desde la confluencia. En magenta, otra capa, que acompaña al río La Caldera a la altura del pueblo. En naranja, el parcelario urbano; en azul, el rural. El visor no declara sobre qué norma se trazó el límite.}%
{Captura del visor \textit{mapas.salta.gob.ar}, 15 de septiembre de 2026. Se recortó el margen
superior con el cuadro de coordenadas de la consulta; no se alteró el mapa.}

\lamina{lam-idesa-municipios.jpg}{lam:idesamun}%
{Límites municipales y departamentales en la exportación del visor de IDESA}%
{\textbf{Exportación «Mapas IDESA»: capas «Límites Municipales» (punteado verde) y «Departamento» (punteado oliva).} Dentro del departamento La Caldera, el límite entre los dos municipios recorre el oeste en dirección noroeste--sudeste y alcanza el corredor de la Ruta 9 unos kilómetros al sur del pueblo, marcado con el alfiler; Vaqueros queda al sudoeste de esa línea. Los demás trazos verdes corresponden a límites de municipios vecinos.}%
{Exportación en PDF del visor de IDESA (Leaflet, Instituto Geográfico Nacional y OpenStreetMap), 16 de septiembre de 2026. Se recortó la zona del departamento y se reubicaron sobre el mapa, sin alterarlos, la leyenda y el indicador de norte; la barra de escala se redibujó a la misma medida que la original.}

\section{Un instrumento serio}

La hipótesis intuitiva sobre un municipio periurbano desbordado supone un gobierno local sin instrumentos. En La Caldera esa hipótesis es falsa.

\begin{ficha}{Código de Ordenamiento Territorial del Municipio de La Caldera}
\begin{itemize}[leftmargin=*,nosep]
\item \textbf{Norma marco:} el PDUA aprobado por Ordenanza Nº 643 del 25/11/2015, que este Código reglamenta.
\item \textbf{Estructura:} 9 títulos, 152 artículos, 8 anexos, atlas de 8 planos que el Código declara confeccionados sobre capas georreferenciadas y que publica sólo como imagen.
\item \textbf{Ámbito:} todo el departamento excepto el municipio de Vaqueros (leyes 2131/47 y 4987/74).
\item \textbf{Marco superior:} Ley provincial 7543 de bosques nativos y Ley 7107 del Sistema Provincial de Áreas Protegidas.
\item \textbf{Autoridad de aplicación:} Secretaría de Obras Públicas, hasta tanto se cree la Secretaría de Planeamiento (art. 23).
\item \textbf{Órgano Técnico de Aplicación:} Obras Públicas + Medio Ambiente + Tierra y Hábitat. Dictamina de manera ineludible sobre emprendimientos de impacto, Planes de Sector, Proyectos Especiales, convenios urbanísticos y usos no conformes.
\item \textbf{Ejecutivo al momento de la redacción:} intendente Daniel Alejandro Escalera\index{Escalera, Daniel Alejandro}; Medio Ambiente, Diego León Vieira Nobre; Obras Públicas, Ing. Luis Borelli; Tierra y Hábitat, Ing. Daniel Cantone; topógrafo Néstor Saavedra.
\item \textbf{Concejo Deliberante:} presidente Lucio Teodoro Sarapura; vicepresidentes Alfio Edgardo Sumbay\index{Sumbay, Alfio Edgardo} y César Dionisio Sumbay; secretaria María Susana Campero.
\item \textbf{Autoría técnica:} Arq. Mariana Zoricich, Dra. Alejandra Cau Cattan, Diego León Vieira Nobre. Colaboración: Arq. Adriana Krumpholz (Provincia), INTA Valle de Lerma.
\end{itemize}
\end{ficha}

\section{Cuatro artículos que ordenan el libro}

\textbf{Artículo 42, inciso 1.} Al definir el Área Urbana, el Código establece que ésta se compone del área urbana primaria y de los \textbf{núcleos urbanos espontáneos o no planificados y de menor consolidación}, vinculados por la Ruta Nacional 9, a los cuales la administración brindará como máximo alumbrado público y recolección de residuos en vías principales. Y los nombra: hacia el sudoeste, \textbf{El Nogalar, El Durazno y Potreros de La Caldera}; hacia el este, Valle Alegre, La Calderilla y La Misión. El inciso 2 agrega que el ordenamiento ``desalienta de manera rotunda la proliferación de estos núcleos dispersos'' y tiende a limitar su extensión y a restringir el desarrollo de nuevos.

\textbf{Zona PU.} Parque Urbano de ribera, sector \textbf{no urbanizable} destinado a uso recreativo, protección ambiental y paisaje, sobre \textbf{ambas márgenes del río La Caldera}, como eje verde del sistema de infraestructura verde municipal. Existe además la Zona IB, para sectores preexistentes de industrias básicas de explotación de áridos, calificadas como actividades incómodas o incompatibles con las residenciales.

\textbf{Artículo 30.} La tributación debe estimular la densificación y el completamiento de las áreas urbanas y la preservación de las rurales y naturales, y dispone expresamente: ``se deberá gravar de manera permanente lotes no edificados en áreas urbanas y loteos o urbanizaciones en áreas no indicadas para tal fin''. Los montos serán proporcionales a superficie, potencial constructivo e infraestructura servida.

\textbf{Artículos 135 a 149.} El Título 7, Capítulo IV establece el procedimiento completo de urbanización o loteo bajo la figura de Plan de Sector: certificado de uso conforme, certificado de prefactibilidad, aprobación de prefactibilidad de anteproyecto, Certificado de Aptitud Ambiental Municipal, Certificado de Aptitud Técnica, permiso de inicio de obra, \textbf{garantía de ejecución de obras de urbanización} (art. 144), \textbf{licencia de comercialización} (art. 145), control de avance, certificado de final de obra, solicitud de matrícula y un artículo final sobre acceso a la información y defensa del consumidor. El Capítulo III regula el Procedimiento de Evaluación de Impacto Ambiental y Social en dieciséis artículos, con categorización, estudio, dictamen técnico, audiencia pública y certificado.

\begin{cautela}
Los artículos 144 y 145 son la clave operativa del libro. El Código exige garantía de ejecución de las obras \textit{antes} de la comercialización y somete la venta a licencia municipal expresa. Para los loteos hoy inundados sólo hay dos posibilidades: si esos actos existen, el municipio los otorgó sobre obras que no se ejecutaron; si no existen, la comercialización fue ilegal y hubo compradores de buena fe perjudicados. No hay una tercera. Ambas ramas son documentables y ninguna fue verificada por este relevamiento.

El artículo 42 tiene además una consecuencia procesal directa. El municipio clasificó por ordenanza, como no planificados y de servicios mínimos, a los tres barrios que tres años después serían la parte actora del amparo colectivo por los aluviones. No puede alegar desconocimiento del carácter irregular de esos núcleos, porque lo declaró él mismo en una norma de orden público. Ese artículo debería incorporarse como prueba documental en la etapa de ejecución de la sentencia.
\end{cautela}

\section{La distancia entre la norma y el territorio}

El artículo 30 explica una declaración pública del intendente. En enero de 2026, ante las inundaciones, sostuvo que las propiedades de esos barrios no pagan impuestos municipales. El Código, sancionado por su propio Concejo, ordena gravar de manera permanente precisamente a los loteos en áreas no indicadas para tal fin. La declaración, leída contra el artículo, no describe una injusticia que el municipio padece: describe un artículo de su propio Código que no está aplicando. El capítulo \ref{cap:hacienda} muestra por qué.

La Zona PU también resuelve, en el papel, un conflicto que la realidad no resolvió: la protección del monte ribereño que el centro vecinal de La Calderilla reclamaba en enero de 2026 estaba vigente desde la sanción del Código, y la extracción de áridos sobre el lecho, condenada por el amparo de 2021, tenía zona propia y calificación de incompatibilidad.

El problema del departamento no es la ausencia de norma sino la distancia entre la norma y el territorio. Eso desplaza el objeto del libro: la pregunta ya no es qué le falta al municipio, sino por qué un municipio que se dotó de un instrumento de esta calidad produce, sobre el mismo suelo que el instrumento clasifica, los resultados que el instrumento prohíbe.

El libro ofrece dos respuestas, y las dos son necesarias. La primera es procedimental y está en este capítulo: el Código contiene su propia vía de excepción. La segunda es material y está en el capítulo \ref{cap:hacienda}: el municipio no tiene con qué aplicarlo.

\section{Qué se aprobó bajo este Código, y qué no}

El Código es de 2016. Los emprendimientos que este libro documenta en el municipio de La
Caldera se aprobaron, en su mayoría, por resolución del \textbf{Ministerio provincial} y no
por ordenanza municipal: el Barrio La Misión y Chacras del Gallinato en 2015, El Durazno en
2021.

Ese dato importa más de lo que parece. Un Código municipal de ciento cincuenta y dos
artículos, con nueve títulos, ocho anexos y atlas georreferenciado, coexiste con un
procedimiento provincial de aprobación de fraccionamientos que no lo menciona. La Resolución
646 D/21 que aprobó el plano de El Durazno se funda en la Ley 2308 y en el Decreto 1682/19,
y sus condiciones para la asignación de matrículas son dos: libre deuda y donación del
espacio institucional. \textbf{Ninguno de los diez actos que este Código exige ---enumerados en el apéndice \ref{cap:matriz}--- aparece en ella}, y eso que el propio anexo del decreto provincial manda iniciar el trámite en el municipio con certificado de factibilidad de localización, pre-visado del plano y certificado de aptitud ambiental (capítulo \ref{cap:loteo}).

Este libro no puede afirmar que esos certificados no existan: el capítulo \ref{cap:loteo} documenta que el
expediente tiene fojas que no se consultaron, y el municipio intervino en él. Afirma que
\textbf{no constan en el acto que aprueba el plano}, que es el único documento público del
trámite, y que la matriz del apéndice \ref{cap:matriz} publica por eso una columna entera de casilleros
vacíos.

\section{La excepción que se tramita}

En diciembre de 2020, tras denuncia formulada por la propia Municipalidad de La Caldera ante la Secretaría de Medio Ambiente provincial, el organismo dispuso la clausura preventiva del loteo privado \textit{La Falda}, en La Calderilla, por movimientos de suelo, apertura de caminos y depredación ambiental, con incumplimiento de la ley provincial 7070 y sin autorización municipal. Los trabajos se realizaban sobre la matrícula 154, entre los kilómetros 16 y 17,5 de la Ruta Nacional 9, sobre ladera de cerro; la apertura de caminos y demarcación estuvo a cargo de la empresa Norte Áridos. Esa misma empresa era contratista de la Provincia en el mismo paraje: en octubre de 2024 el Ministerio de Infraestructura aprobó el convenio de rescisión que celebró con «la contratista Norte Áridos S.R.L.» para la obra «Optimización del servicio de La Calderilla --- nueva red distribuidora» (Resolución 154/24, B.O.\ Nº 21.828).

Según el comunicado del propio municipio, el titular había hecho presentaciones para obtener la factibilidad, pero el personal técnico determinó que \textbf{incumple el Código porque éste no permite la urbanización en el sector, ya que degrada el medioambiente}. Y a continuación se le indicó que presentara ``un proyecto especial que se ajuste al COTM y cuente con el Certificado de Aptitud Ambiental otorgado por Medioambiente de la provincia, para que eleve al Concejo Deliberante, lo analice y apruebe''.

Léase la secuencia. El municipio establece que el sector no es urbanizable según su propio Código, porque la urbanización degradaría el ambiente. Y a continuación indica la vía por la cual podría urbanizarse de todos modos: la figura del Proyecto Especial del artículo 31, que somete el caso a dictamen del Órgano Técnico y aprobación por ordenanza.

No hay irregularidad. El Proyecto Especial existe justamente para desarrollos que requieren tratamiento particularizado, y el municipio actuó denunciando y clausurando. Pero el efecto combinado es que \textbf{la calificación de no urbanizable no opera como prohibición sino como requisito de procedimiento agravado}. Lo que en el plano de zonificación es un límite, en el trámite es una puerta con más pasos.

Ésa es la primera respuesta a la pregunta del apartado anterior. No es que el Código se incumpla: es que contiene su propia vía de excepción, y la excepción se tramita.

Cinco años después, en enero de 2026, el presidente del centro vecinal de La Calderilla declaró a una radio provincial que había numerosas ofertas de loteos sobre los cuales el municipio no respondía si estaban aprobados; que ciertos emprendimientos desmontaron el lote completo incluida la franja de monte ribereño sin autorización de la Secretaría de Medio Ambiente; y que, pese a que el emprendimiento habría sido multado por el organismo ambiental, el municipio avanzó con una audiencia pública para tratar el loteo, en la cual la presidenta del Concejo Deliberante agradeció la inversión mientras el proyecto continuaba penalizado. Para esa audiencia, el buscador de avisos del Boletín Oficial no devuelve convocatoria entre 2019 y 2026 ---fuera de los tramos leídos edición por edición, eso sólo significa que el barrido no la encontró---: la única audiencia sobre un loteo del departamento que devuelve en ese lapso es la de Lares de la Inmaculada, en Vaqueros; la otra convocatoria del período, la de Vialidad por la Ruta 9 en noviembre de 2023, es de una obra vial (capítulo \ref{cap:vaqueros})---, y desde enero de 2025 el Concejo lo presiden hombres; si el relato es exacto, la audiencia es anterior, de cuando lo presidían Yamila Sastre o Ana Sumbay (capítulo \ref{cap:infraestructura}).

Si el relato es exacto --- y proviene de un dirigente vecinal identificado, en declaración pública --- describe el punto donde el circuito falla: no en la ausencia de control, sino en la desconexión entre el nivel que sanciona y el que habilita. La provincia multa; el municipio agradece la inversión. Ninguno de los dos actos es ilegal por separado. La combinación produce un resultado que ningún funcionario decidió y que todos ejecutaron.

Ésa es, en escala municipal, la definición operativa de dispositivo que el libro anterior sostuvo durante cuatrocientas páginas.

